\documentclass[a4paper, 12pt]{article}

\usepackage{utils}

\renewcommand*{\today}{24 March 2026}

\begin{document}

\hotbox{Algèbre bilinéaire}{CM 8}{\today}

\subsection{Projections et symétries orthogonales}

\begin{definition}
    Soit $(E, \langle \cdot, \cdot \rangle)$ un espace vectoriel euclidien de dimension finie et $V \subset E$ un sev.

    Soit $(v_1, \dots, v_k)$ une base orthonormée de $V$. Pour tout $x \in E$, on définit la \textbf{projection orthogonale} de $x$ sur $V$ par
    $$
    p(x) = \sum_{i=1}^k \langle x, v_i \rangle v_i
    $$
\end{definition}

\begin{proposition}{}{}
    Soit $(E, \langle \cdot, \cdot \rangle)$ un espace vectoriel euclidien de dimension finie et $V \subset E$ un sev.
    
    Pour tout $x \in E$, on a
    \begin{enumerate}
        \item $p(x) - x \in V^\perp$
        \item $\norm{p(x) - x} = min \{\norm{y - x} \mid y \in V\}$
        \item $p(x)$ ne dépend pas du choix de la base orthonormée de $V$
    \end{enumerate}
\end{proposition}

\begin{demonstration}
    \begin{enumerate}
        \item .
        \item .
        \item Soit $(u_1, \dots, u_k)$ une autre base orthonormée de $V$.
        $$
        p'(x) = \sum_{i=1}^k \langle x, u_i \rangle u_i
        $$
        pour tout $x \in E$.

        Ainsi, on a
        $$
        p(x) - p'(x) = p(x) - x - (p'(x) - x) \in V^\perp
        $$
        et
        $$
        p(x) - p'(x) \in V
        $$
        Donc $p(x) - p'(x) = \vec{0}$ et $p(x) = p'(x)$
    \end{enumerate}
\end{demonstration}

\begin{proposition}{}{}
    Soit $(E, \langle \cdot, \cdot \rangle)$ un espace vectoriel euclidien de dimension finie et $V \subset E$ un sev.
    
    $$
    Im(p) = V \quad \text{et} \quad Ker(p) = V^\perp
    $$
\end{proposition}

\begin{remarque}
    Si on décompose $x$ en $x = x' + x''$ avec $x' \in V$ et $x'' \in V^\perp$, alors $p(x) = x'$
\end{remarque}

\begin{definition}
    Soient $(E, \langle \cdot, \cdot \rangle)$ un espace euclidien de dimension finie et $V \subset E$ un sev.

    On définit la \textbf{symétrie orthogonale} de $x$ sur $V$ par
    $$
    s(x) = x' - x''
    $$
    où $x = x' + x''$ avec $x' \in V$ et $x'' \in V^\perp$
\end{definition}

\begin{proposition}{}{}
    Soit $(E, \langle \cdot, \cdot \rangle)$ un espace euclidien avec de dimension finie et $V \subset E$ un sev.

    \begin{enumerate}
        \item $s = 2p - id_E$
        \item $s^2 = id_E$
        \item $Ker(s - id_E) = V$ et $Ker(s + id_E) = V^\perp$
    \end{enumerate}
\end{proposition}

\begin{demonstration}
    \begin{enumerate}
        \item Par définition de $s$, on a pour tout $x \in E$,
        $$
        s(x) = p(x) - (x - p(x)) = 2p(x) - x
        $$
        donc $s = 2p - id_E$
        \item Laissé au lecteur.
        \item On a
        $$
        Ker(s - id_E) = Ker(2p - 2id_E) = Ker(p - id_E) = \{x \in E \mid p(x) = x\} = V
        $$
        et
        $$
        Ker(s + id_E) = Ker(2p) = Ker(p) = \{x \in E \mid p(x) = \vec{0}\} = V^\perp
        $$
    \end{enumerate}
\end{demonstration}

\begin{proposition}{}{}
    Soit $\sigma: E \to E$ une application linéaire telle que
    \begin{enumerate}
        \item $\sigma^2 = id_E$
        \item $Ker(\sigma - id_E) = Ker(\sigma + id_E)^\perp$
    \end{enumerate}
    Alors $\sigma$ est une symétrie orthogonale par rapport à $Ker(\sigma - id_E)$
\end{proposition}

\subsection{Orientation, produit mixte et produit vectoriel}

\begin{definition}
    Soit $(E, \langle \cdot, \cdot \rangle)$ un espace vectoriel euclidien de dimension finie $n$.

    Soit $(x_1, \dots, x_p) \in E^p$, avec $p \leq n$, on pose
    $$
    G(x_1, \dots, x_p) = (\langle x_i, x_j \rangle)_{1 \leq i, j \leq p} \in M_{pp}(\R)
    $$
    $G$ est appelée la \textbf{matrice de Gram} de $(x_1, \dots, x_p)$ et son déterminant est appelé le \textbf{déterminant de Gram} de $(x_1, \dots, x_p)$
    $$
    gram(x_1, \dots, x_p) = det(G(x_1, \dots, x_p))
    $$
\end{definition}

\begin{proposition}{}{}
    Soit $(E, \langle \cdot, \cdot \rangle)$ un espace vectoriel euclidien de dimension finie $n$ et $(x_1, \dots, x_p) \in E^p$ avec $p \leq n$.

    Notons $C_i$ la $i$-ème colonne de $G(x_1, \dots, x_p)$, alors

    \begin{enumerate}
        \item Pour tout $\lambda_1, \dots, \lambda_p \in \R$,
        $$
        \sum_{i=1}^p \lambda_i x_i = 0 \iff \sum_{i=1}^p \lambda_i C_i = 0
        $$

        \item
        $(x_1, \dots, x_p)$ est libre si et seulement si $gram(x_1, \dots, x_p) \neq 0$.
    \end{enumerate}
\end{proposition}

\begin{demonstration}
    \begin{enumerate}
        \item
        \begin{description}
            \item[Sens direct.] Supposons que $\sum_{i=1}^p \lambda_i x_i = 0$. 
            On a
            $$
            \sum_{i=1}^p \lambda_i C_i = \sum_{i=1}^p \lambda_i (\langle x_1, x_i \rangle, \dots, \langle x_p, x_i \rangle)^t = (\langle \sum_{i=1}^p \lambda_i x_i, x_1 \rangle, \dots, \langle \sum_{i=1}^p \lambda_i x_i, x_p \rangle)^t = 0
            $$ 
            \item[Sens réciproque.] Supposons que $\sum_{i=1}^p \lambda_i C_i = 0$.
            
            C'est à dire que pour tout $1 \leq j \leq p$,
            $$
            \sum_{i=1}^p \lambda_i \langle x_j, x_i \rangle = 0 \implies \langle x_j, \sum_{i=1}^p \lambda_i x_i \rangle = 0 \\
            $$
            On peut poser $v = \sum\limits_{i=1}^p \lambda_i x_i$. Ainsi $v$ est orthogonal à tous les $x_j$ pour $1 \leq j \leq p$. Donc $v$ est orthogonal à tous les éléments de $Vect(x_1, \dots, x_p)$.
            En particulier, $v$ est orthogonal à lui même, d'où $v = 0$ (car le produit scalaire est défini) et $\sum_{i=1}^p \lambda_i x_i = 0$.
        \end{description}
    \end{enumerate}
\end{demonstration}

\subsection{Orientation}

\begin{definition}
    Soit $(E, \langle \cdot, \cdot \rangle)$ un espace vectoriel euclidien de dimension finie $n$.

    Soient $\mathcal{B}$ et $\mathcal{B}'$ deux bases de $E$.
    
    On dira que $\mathcal{B}$ et $\mathcal{B}'$ ont la même orientation (notée $\mathcal{B} \sim \mathcal{B}'$) si $det(P_{\mathcal{B} \to \mathcal{B}'}) > 0$.

    En partant d'une base $\mathcal{B}$ de référence de $E$, les bases de même orientation que$\mathcal{B}$ seront dits \textbf{directes} et les autres \textbf{indirecte}.
\end{definition}

\begin{remarque}
    L'orientation d'un espace vectoriel euclidien de dimension finie $n$ est une relation d'équivalence qui divise $E$ en deux classes d'équivalence.
\end{remarque}

\begin{remarque}
    Toutes surfaces n'est pas obligatoirement orientable. Par exemple, la bouteille de Klein ou le ruban de Möbius.
\end{remarque}

\begin{proposition}{}{}
    Soit $(E, \langle \cdot, \cdot \rangle)$ un espace vectoriel euclidien de dimension finie $n$.

    Soient $(e_1, \dots, e_n)$ une base orthonormée de $E$ et $(x_1, \dots, x_n) \in E^n$. On a pour $1 \leq i \leq n$
    $$
    x_i = \sum_{j=1}^n \lambda_{ij} e_j
    $$
    On pose $\Lambda = (\lambda_{ij})_{1 \leq i, j \leq n}$ alors $\Lambda(^t \Lambda) = G(x_1, \dots, x_n)$
\end{proposition}

\begin{demonstration}
    On a
    $$
    (\Lambda(^t \Lambda))_{ij} = \sum_{k=1}^n \lambda_{ik} \lambda_{jk} = \langle x_i, x_j \rangle = G(x_1, \dots, x_n)_{ij}
    $$
    Donc $\Lambda(^t \Lambda) = G(x_1, \dots, x_n)$
\end{demonstration}

\begin{definition}[(Produit mixte)]
    Soit $(E, \langle \cdot, \cdot \rangle)$ un espace vectoriel euclidien de dimension finie $n$.
    
    Soit $\mathcal{B}$ une base orthonormée directe de $E$.
    
    $det_\mathcal{B}(x_1, \dots, x_n)$ est appelé le \textbf{produit mixte} de $(x_1, \dots, x_n)$ et est noté $[x_1, \dots, x_n]$.
\end{definition}

\begin{proposition}{}{}
    Le produit mixte ne dépend pas du choix de la base orthonormée directe de $E$.
\end{proposition}

\begin{demonstration}
    Soient $\mathcal{B}$ et $\mathcal{B}'$ deux bases orthonormées directes de $E$.

    Alors
    $$
    det_\mathcal{B}(x_1, \dots, x_n) = det(P_{\mathcal{B} \to \mathcal{B}'}) det_\mathcal{B'}(x_1, \dots, x_n) = det_\mathcal{B'}(x_1, \dots, x_n)
    $$
    car $det(P_{\mathcal{B} \to \mathcal{B}'}) = 1$.

    D'où le résultat.
\end{demonstration}

\begin{remarque}
    Le produit mixte correspond géométriquement au volume orienté du parallélépipède formé par les vecteurs $x_1, \dots, x_n$.
\end{remarque}

\begin{proposition}{}{}
    Soit $(E, \langle \cdot, \cdot \rangle)$ un espace vectoriel euclidien de dimension finie $n$ et $(x_1, \dots, x_n) \in E^n$.

    $$
    [x_1, \dots, x_n]^2 = gram(x_1, \dots, x_n)
    $$
\end{proposition}

\begin{theoreme}{Produit vectoriel}{}
    Soit $(E, \langle \cdot, \cdot \rangle)$ un espace vectoriel euclidien de dimension finie $n$.
    Pour toute famille $(x_1, \dots, x_{n-1}) \in E^{n-1}$, il existe un unique vecteur $a$ (noté $x_1 \wedge \dots \wedge x_{n-1}$) tel que pour tout $x \in E$,
    $$
    [x_1, \dots, x_{n-1}, x] = \langle a, x \rangle
    $$

    Ce vecteur $a$ est appelé le \textbf{produit vectoriel} de $(x_1, \dots, x_{n-1})$.
\end{theoreme}

\begin{demonstration}
    Fixons $(x_1, \dots, x_{n-1}) \in E^{n-1}$.
    \begin{description}
        \item[Unicité]
        On suppose qu'il existe $a'$ tel que pour tout $x \in E$,
        $$
        [x_1, \dots, x_{n-1}, x] = \langle a', x \rangle
        $$
        Ainsi, pour tout $x \in E$,
        $$
        \langle a, x \rangle = \langle a', x \rangle \iff \langle a - a', x \rangle = 0
        $$
        Ainsi $a - a'$ est orthogonal à tous les éléments de $E$, en particulier à lui même, d'où $a - a' = \vec{0}$ et 
        d'où $a = a'$

        \item[Existence (1ère méthode)]
        On considère
        $$
        \funcdef{\varphi}{E}{\R}{x}{[x_1, \dots, x_{n-1}, x]}
        $$
        C'est une forme linéaire sur $E$.

        Soit $\mathcal{B} = (e_1, \dots, e_n)$ une base orthonormée de $E$.

        Alors les formes $\funcdef{e_i^*}{E}{\R}{x}{\langle e_i, x \rangle}$ forment une base de $E^*$.

        Ainsi $\varphi = \sum_{i=1}^n \alpha_i e_i^* \in E^*$ pour certains $\alpha_i \in \R$.

        Maintenant posons
        $$
        a = \sum_{i=1}^n \alpha_i e_i \in E
        $$
        On a
        $$
        [x_1, \dots, x_{n-1}, x] = \varphi(x) = \sum_{i=1}^n \alpha_i e_i^*(x) = \sum_{i=1}^n \alpha_i \langle e_i, x \rangle = \langle a, x \rangle
        $$
        \item[Existence (2ème méthode)]
        On change les notations et on pose $x_n = x$ et $a_n = a$.
        On considère $\mathcal{B} = (e_1, \dots, e_n)$ la base canonique de $E$ (qui est orthonormée).
        On a
        \begin{align*}
            [x_1, \dots, x_{n-1}, x_n] &= det_\mathcal{B}(x_1, \dots, x_{n-1}, x_n) \\
            &= \sum_{\sigma \in S_n} \varepsilon(\sigma) \prod_{i=1}^n \lambda_{i \sigma(i)} \\
            &= \sum_{i=1}^n \lambda_{i, n} \sum_{\sigma \in \Lambda_{i, n}} \varepsilon(\sigma) \prod_{j=1, j \neq i}^n \lambda_{j \sigma(j)} \quad \text{avec} \quad \Lambda_{i, n} = \{ \sigma \in S_n \mid \sigma(i) = n \} \\
            &= \sum_{i=1}^n \lambda_{i, n} a_{i,n} \quad \text{avec} \quad a_{i, n} = \sum_{\sigma \in \Lambda_{i, n}} \varepsilon(\sigma) \prod_{j=1, j \neq i}^n \lambda_{j \sigma(j)} \\
            &= \langle x_n, a_n \rangle
        \end{align*}
        On remarque que les $a_{i, n}$ sont les cofacteurs.
    \end{description}
\end{demonstration}

\begin{remarque}
    Le produit vectoriel correspond géométriquement à un vecteur orthogonal à tous les $x_i$ pour $1 \leq i \leq n-1$ et dont la norme est égale à l'aire du parallélogramme formé par les $x_i$.
\end{remarque}

\end{document}